\section{Conclusion, Limitations and Future Work}
\label{sec:cfw}

\subsection{Conclusion}

This project develops an integrated relational database for selected
environmental records from several Bangladeshi organisations. The final work
brings the approved source data into a consistent relational structure while
preserving differences in record grain, source coverage and missing
information.

The project also establishes a clear connection between the source material,
the Entity--Relationship model, the normalized relations and the implemented
SQLite database. Data are included only when their meaning, identifiers and
relationships can be supported by the available evidence. Where information
is missing or uncertain, the database preserves that limitation rather than
introducing unsupported values.

Within the approved scope, the resulting database provides a structured and
reproducible basis for storing and querying the selected environmental data.
It contains \mDbTables{} tables, \mDbViews{} views, \mDbColumns{} table columns
and \mDbRows{} rows covering \mYearMin--\mYearMax. It can also serve as a
starting point for later groups that may continue the environment-sector work
in future academic sessions.

\subsection{Limitations and Future Work}
\label{subsec:future}

The current implementation completes the approved scope of this project, but
some limitations remain because of the available source material and the
present deployment environment. These are not unfinished tasks of the current
team. Instead, they identify areas that later groups may continue if the
project is extended. Table~\ref{tab:limitations-future-work} separates each
current limitation from its present handling and possible future continuation.

\begin{longtable}{L{0.20\textwidth}L{0.32\textwidth}L{0.32\textwidth}}
\caption{Current limitations and directions for future continuation}
\label{tab:limitations-future-work}\\
\toprule
\textbf{Current Limitation} & \textbf{Current Handling} &
\textbf{Future Continuation} \\
\midrule
\endfirsthead
\multicolumn{3}{l}{\small\textit{Table~\thetable{} continued}}\\
\toprule
\textbf{Current Limitation} & \textbf{Current Handling} &
\textbf{Future Continuation} \\
\midrule
\endhead
\midrule
\multicolumn{3}{r}{\small\textit{Continued on the next page}}\\
\endfoot
\bottomrule
\endlastfoot
Uneven data coverage &
Different sources provide records for different locations, variables, time
periods and observation grains. Only the available source records are retained. &
New environmental datasets can be added when reliable sources, clear record
grain and suitable identifiers are available. \\
Unconnected candidate datasets &
Some reviewed datasets cannot be connected reliably with the approved model
because they lack dependable joining fields or suitable structure. &
These datasets can be reconsidered if improved versions or reliable reference
keys become available. \\
Incomplete geographic mapping &
\mUnmappedStations{} of the 56 climate stations do not have a published
district mapping in the retained sources. Their district value remains null. &
An official geographic reference source can be added after the
station--district relationships are verified. \\
Source inconsistencies &
Some values, headings and names remain uncertain when the source does not
provide enough evidence for a definite correction. &
Newer publications or additional official sources may resolve these cases when
stronger evidence becomes available. \\
Manual review requirements &
Irregular spreadsheet layouts and source-specific structures still require
inspection during extraction and verification. &
Extraction tests can be extended for duplicate records, naming differences,
missing-value patterns, unexpected values and publisher layout changes. \\
Narrow industry coverage &
The database retains one industry category and only \mIndustryRows{} fiscal
rows, so it cannot represent all industrial wastewater activity. &
Additional industry data can be included only when its categories, periods and
measurement meanings can be matched reliably. \\
No peer database for comparison &
The peer group did not provide an executable SQLite database. The reported
competency evaluation therefore covers Group-07's implementation only. &
No peer comparison is required for the present report. A later verified
database could be tested with the same competency questions as separate work. \\
Single-user SQLite deployment &
SQLite is suitable for the current portable course-project implementation but
is not intended here for simultaneous authenticated multi-user access. &
The normalized schema can later be migrated to a server database if multi-user
access and centralized management are required. \\
Direct query-based access &
Users currently access the stored records mainly through SQL queries and
database tools. &
A simple web interface, dashboard or Application Programming Interface (API)
can be built over the existing database. \\
\end{longtable}

\subsection{Verification Completed}
\label{subsec:verification-completed}

The approved final ERD, schema definition, normalized CSV headers and SQLite
database are checked against the same relation and key definitions. Automated
tests rebuild the database, run integrity and foreign-key checks, execute the
competency queries and verify backup and restoration by comparing the restored
database contents. The current database reports \mDbIntegrity{} for
\texttt{PRAGMA integrity\_check} and \mDbFkViolations{} foreign-key violations.
These checks are already implemented and are not being deferred as future work.

If the project is continued by a future batch, the present database, final ERD,
normalized structure, source records and verification files should be used as
the starting point rather than rebuilding the environment-sector database from
the beginning. Any new dataset should first be checked for source reliability,
record grain, identifiers and compatibility with the existing model.

Future work should therefore focus mainly on extending the database where new
evidence becomes available and improving the supporting tools around the
existing structure. The extraction tests should also be updated whenever a
publisher changes a workbook layout. This approach allows later groups to
continue the project while preserving the consistency and source traceability
established through the retained normalization and review files.
